\section{模型的评价与推广}

\subsection{模型优点}

首先，本文从有限视线段出发保留了真目标完整圆柱几何，而不是只检验圆柱中心点。该判据直接对应“烟幕遮住完整目标投影”的物理含义，能够避免中心线通过烟幕但目标边缘仍暴露的乐观估计。其次，所有问题共用同一组导弹、无人机、烟幕弹和云团运动方程，后续问题只增加决策规模与协同关系，保证五个问题之间具有清晰的继承链。

再次，本文用阈值事件求根确定遮蔽区间端点，再按集合并集计算多弹时长。该处理同时解决了固定时间步长误差和重叠区间重复计数问题。对于问题五，先进行无人机--导弹效能筛查，再交替优化任务分配和连续参数，显著缩小了高维混合搜索空间；最终方案只使用具有正收益或明显边界裕量的烟幕弹，便于实际执行。

最后，求解过程具有较强可复现性。各问均输出完整决策变量、投放点、起爆点、有效区间和高密度复核结果；问题三至五按官方模板生成 Excel 文件，图表直接读取同一结果数据，避免论文、代码和附件之间出现手工转录误差。

\subsection{模型局限}

本文按题面忽略风场、空气阻力、烟幕扩散与浓度衰减，并把烟幕有效区固定为半径 $10\,\mathrm{m}$ 的球体。真实环境中风速和湍流会使云团发生平移、拉伸和非均匀稀释，固定球模型可能低估或高估某些方向的遮蔽时长。无人机瞬时转向和始终匀速平飞也是理想化假设，没有考虑转弯半径、加速度限制和飞行安全间隔。

完整圆柱判据通过边界离散近似连续目标，尽管 $5760$ 点终算已显示良好收敛，但仍不是连续几何的解析证明。差分进化及分层任务分配均属于启发式方法，能提供高质量可行解，却不能保证在全部离散分配和连续参数组合中达到严格全局最优。问题五的目标采用三枚导弹遮蔽时长之和，默认三种威胁权重相同；若实际任务更重视最危险导弹或要求三者最低保护水平，则需修改目标函数。

\subsection{稳健性改进方向}

若获得风速、风向和烟幕浓度数据，可把球心运动扩展为带风漂移的随机微分模型，并将有效半径替换为随时间变化的等浓度面。对投放时刻、速度和引信延时引入误差变量后，可用情景法或机会约束最大化期望遮蔽时长，同时限制失效概率。工程应用中还可把目标改为
\begin{equation}
 \max\left\{\sum_{j=1}^{3}w_jT_j-\lambda C,
 \quad \min_{j=1,2,3}T_j\right\},
\end{equation}
其中 $w_j$ 表示威胁权重，$C$ 表示弹药和机动代价，$\lambda$ 为成本系数。第一种形式适合总体效益最大化，第二种形式适合强调最弱防护对象的公平配置。

\subsection{推广应用}

本文的核心并不限于烟幕弹。只要防护装置能够形成随时间运动的空间作用域，就可以把“作用域与观测视线相交”作为通用覆盖判据。例如，无人机通信中继可把烟幕球替换为有效通信域，目标函数改为链路连通时间；森林灭火可把导弹视线替换为火线传播路径，把烟幕弹替换为空投阻燃剂；移动传感器部署则可把完整圆柱替换为待观测区域，以覆盖时间并集评价多平台协同效果。多目标场景中的“配对探针--任务分配--连续细化--统一复核”框架也适用于带资源容量和时空窗口的应急调度问题。

\endinput
\section{模型的评价与推广}
